% Part: proof-theory
% Chapter: cut-elimination
% Section: midsequent

\documentclass[../../../include/open-logic-section]{subfiles}

\begin{document}

\olfileid[pt]{pt}{cut}{mid}

\olsection{O Teorema do Sequente Médio e o Teorema de Herbrand}

Uma consequência importante do teorema da eliminação do corte é o
teorema do sequente médio. Ele afirma que, se existe !!a{proof}~$\pi$
de~$\Gamma \Sequent \Delta$ na qual toda !!{formula} é prenexa, então
existe !!a{proof}~$\pi'$ que contém um sequente $S \ident \Pi \Sequent
\Lambda$ (chamado de \emph{sequente médio}) tal que todas as inferências acima de
$S$ em~$\pi'$ são proposicionais e todas as inferências abaixo de~$S$ são
inferências de quantificadores. (!!^a{formula}~$!A$ é \emph{prenexa} (ou está na
\emph{forma prenexa}) se tem a forma
\[\mathsf{Q}_1x_1\dots\mathsf{Q}_nx_n\,!B(x_1,\dots,x_n)\] em que cada
$\mathsf{Q}_i$ é $\lforall$ ou~$\lexists$.) Uma consequência importante
do teorema do sequente médio é o Teorema de Herbrand. Seu
caso geral é um pouco difícil de descrever, mas, em termos muito gerais, tem
a forma seguinte: para toda !!{provable} !!{sentence}~$!A$ da lógica de primeira
ordem, existe uma tautologia proposicional~$!A'$ a partir da qual $!A$ pode ser !!{prove}d. Eis a versão para
!!{sentence}s da forma~$\lexists[x][!B(x)]$, em que $!B$ é
sem quantificadores:

\begin{thm}[Teorema de Herbrand]
Suponha que $!B(x)$ seja sem quantificadores e que $\Proves 
\lexists[x][!B(x)]$. Então existem termos~$t_1$, \dots, $t_n$
tais que $\Proves !B(t_1) \lor \dots \lor !B(t_n)$.
\end{thm}

A !!{formula}~$!B(t_1) \lor \dots \lor !B(t_n)$ é chamada de
\emph{disjunção de Herbrand} de~$\lexists[x][!B(x)]$. Como $!B$ é
sem quantificadores, essa disjunção também é sem quantificadores. Assim, se ela é
!!{provable}, é !!{provable} por uma prova sem corte e, consequentemente,
por !!a{proof} que usa apenas inferências proposicionais. Logo, é
uma tautologia.

A parte difícil do Teorema de Herbrand é mostrar que existe uma disjunção de
Herbrand para toda !!{provable}
!!{formula}~$\lexists[x][!B(x)]$. A recíproca, O.s.\@, que, se
$\lexists[x][!B(x)]$ tem uma disjunção de Herbrand, então é demonstrável,
é fácil. De fato, a disjunção de Herbrand implica $\lexists[x][!B(x)]$.
P.ex., eis !!a{proof} em~$\Log{G3c}$ para o caso em que $n=2$:
\[
\Axiom$!B(t_1) \fCenter \lexists[x][!B(x)], !B(t_1), !B(t_2)$
\Axiom$!B(t_2) \fCenter \lexists[x][!B(x)], !B(t_1), !B(t_2)$
\RightLabel{\LeftR{\lor}}
\BinaryInf$!B(t_1) \lor !B(t_2) \fCenter \lexists[x][!B(x)], !B(t_1), !B(t_2)$
\RightLabel{\RightR{\lexists}}
\UnaryInf$!B(t_1) \lor !B(t_2) \fCenter \lexists[x][!B(x)], !B(t_2)$
\RightLabel{\RightR{\lexists}}
\UnaryInf$!B(t_1) \lor !B(t_2) \fCenter \lexists[x][!B(x)]$
\DisplayProof
\]

Para extrair a disjunção de Herbrand de !!a{proof}~$\pi$ de~$\Sequent
\lexists[x][!B(x)]$, considere o caso em que todas as inferências $\RightR{\lexists}$
em uma !!{proof}~$\pi$ sem corte estão no final:
\[
\AxiomC{}
\RightLabel{$\pi_1$}
\Deduce$ \fCenter \lexists[x][!B(x)], !B(t_1), \dots, !B(t_n)$
\RightLabel{\RightR{\lexists}}
\UnaryInf$ \fCenter \lexists[x][!B(x)], !B(t_2), \dots, !B(t_n)$
\RightLabel{$(n-1) \times \RightR{\lexists}$}
\Deduce$ \fCenter \lexists[x][!B(x)]$
\DisplayProof
\]
Se essas são todas as inferências $\RightR{\lexists}$ em~$\pi$ com
$\lexists[x][!B(x)]$ ativa, então não há outras inferências de quantificadores
em~$\pi_1'$. Isso porque $!B(x)$ não contém outros
quantificadores, e o sequente final não contém quaisquer !!{formula}s à
esquerda. Em outras palavras, $ \fCenter \lexists[x][!B(x)], !B(t_1),
\dots, !B(t_n)$ é o sequente médio de~$\pi$. Além disso, como não há
inferências de quantificadores acima do sequente médio, todos os sequentes acima dele
contêm $\lexists[x][!B(x)]$, e essas ocorrências não são ativas nem
principais. Portanto, elas podem ser apagadas da !!{proof}~$\pi_1$, e obtemos
!!a{proof} de~$\fCenter !B(t_1), \dots, !B(t_n)$ e, usando $n-1$
aplicações de $\RightR{\lor}$, !!a{proof} da disjunção de
Herbrand~$\Sequent !B(t_1) \lor \dots \lor !B(t_n)$.

O problema é que não podemos supor que nem mesmo uma !!{proof} sem corte de
$\Sequent \lexists[x][!B(x)]$ tenha todas as suas inferências \RightR{\lexists}
em um único bloco no final. Pode haver inferências entre essas
inferências \RightR{\lexists} nas quais algum $!B(t_i)$ seja principal.
Daí a necessidade do teorema do sequente médio.

\begin{thm}[Teorema do Sequente Médio]\ollabel{thm:midsequent} Se existe uma
  \Log{G3c}-!!{proof}~$\pi$ de $\Gamma \Sequent \Delta$ em que todas as
  !!{formula}s de $\Gamma$ e $\Delta$ são prenexas, então existe
  !!a{proof}~$\pi'$ que contém um sequente $\Pi \Sequent \Lambda$ tal
  que todas as inferências abaixo dele são inferências de quantificadores e todas
  as inferências acima dele são inferências proposicionais.
\end{thm}

\begin{proof}
  Defina a \emph{ordem} de uma inferência de quantificador em~$\pi$ como o
  número de inferências proposicionais abaixo dela, e a ordem~$o(\pi)$
  de $\pi$ como a soma das ordens de todas as suas inferências de quantificadores.
  Se $o(\pi) = 0$, então nenhuma inferência de quantificador tem uma inferência
  proposicional abaixo dela, e terminamos: como todas as inferências de quantificadores
  têm uma premissa, existe uma única inferência de quantificador mais alta.
  Sua premissa é o sequente médio.

  Se $o(\pi) > 0$, escolha a inferência de quantificador mais baixa entre as de ordem
  $>0$. Como sua ordem é $>0$, deve haver uma inferência proposicional
  abaixo dela. Como ela é a mais baixa, sua conclusão
  não pode ser a premissa de uma inferência de quantificador: caso contrário, essa segunda
  inferência de quantificador, mais baixa, também teria uma inferência proposicional abaixo
  dela. Distinguimos (muitos!) casos segundo qual inferência de quantificador
  temos e qual é a inferência proposicional abaixo dela.
  Como o sequente final consiste apenas em fórmulas prenexas, a
  fórmula principal da inferência de quantificador não pode estar ativa na
  inferência proposicional seguinte. Caso contrário, a fórmula principal
  dessa inferência teria um quantificador no escopo de um
  conectivo proposicional. Pela propriedade da sub-!!{formula}, ela teria
  de ser uma sub-!!{formula} de !!a{formula} no sequente final.
  Assim, a !!{formula} principal da inferência de quantificador é
  !!a{element} do contexto da inferência proposicional inferior.

  Eis dois exemplos:

Primeiro, suponha que tenhamos uma $\RightR{\lforall}$ que termina na premissa
direita de $\LeftR{\lif}$. Então a sub-!!{proof} relevante de~$\pi$
termina na sub-!!{proof}~$\pi_1$.
\[
\AxiomC{}
\RightLabel{$\pi_2$}
\Deduce$\Gamma' \fCenter \Delta', \lforall[x][!A(x)], !B$
\AxiomC{}
\RightLabel{$\pi_3$}
\Deduce$!C, \Gamma' \fCenter \Delta', !A(c)$
\RightLabel{\RightR{\lforall}}
\UnaryInf$!C, \Gamma' \fCenter \Delta', \lforall[x][!A(x)]$
\RightLabel{\LeftR{\lif}}
\BinaryInf$!B \lif !C, \Gamma' \fCenter \Delta', \lforall[x][!A(x)]$
\DisplayProof
\]
Podemos supor que $\pi$ seja regular; assim, a variável própria~$c$ não
ocorre fora de~$\pi_3$; em particular, não ocorre em $!B
\lif !C$ nem em~$\pi_2$. Considere agora a prova~$\pi_1'$:
\[
\AxiomC{}
\RightLabel{$\pi_2'$}
\Deduce$\Gamma' \fCenter \Delta', \lforall[x][!A(x)], !A(c), !B$
\AxiomC{}
\RightLabel{$\pi_3$}
\Deduce$!C, \Gamma' \fCenter \Delta', \lforall[x][!A(x)], !A(c)$
\RightLabel{\LeftR{\lif}}
\BinaryInf$!B \lif !C, \Gamma' \fCenter \Delta', \lforall[x][!A(x)], !A(c)$
\RightLabel{\RightR{\lforall}}
\UnaryInf$!B \lif !C, \Gamma' \fCenter \Delta', \lforall[x][!A(x)], \lforall[x][!A(x)]$
\DisplayProof
\]
em que $\pi_2'$ é obtida de~$\pi_2$ por enfraquecimento com~$!A(c)$ à
direita. (Isso não acrescenta nenhuma inferência, em particular nenhuma
inferência de quantificador que possa afetar a ordem. Também não
viola nenhuma condição de variável própria em~$\pi_2$, pois nenhuma constante
de~$!A(c)$ é usada como variável própria em~$\pi_2$.) A condição de variável própria
de \RightR{\lforall} continua satisfeita, pois $c$ não
ocorre em~$!B \lif !C$.

Invocamos a admissibilidade da contração para
obter !!a{proof}~$\pi_1''$ de $!B \lif !C, \Gamma' \fCenter \Delta',
\lforall[x][!A(x)]$. A demonstração da admissibilidade de
$\RightR{\Contraction}$ para $\lforall[x][!A(x)]$ e a demonstração da
invertibilidade de $\RightR{\lforall}$ usada nessa demonstração não alteraram
a ordem de nenhuma inferência: no máximo, apagaram inferências. Assim,
não há mais inferências de quantificadores em~$\pi_1''$ do que em~$\pi_1$,
e cada inferência de quantificador em~$\pi_1''$ tem o mesmo número de
inferências proposicionais abaixo dela que sua inferência correspondente
em~$\pi_1$---com exceção da última \RightR{\lforall}: ela tinha
uma \LeftR{\lif} abaixo dela em~$\pi$, que movemos para acima dela
em~$\pi_1''$. Substitua $\pi_1$ em $\pi$ por $\pi_1''$. A
!!{proof} resultante tem ordem menor, pois ao menos a ordem da inferência
$\RightR{\lforall}$ foi reduzida em (pelo menos)~$1$.
Agora aplique a hipótese de indução.

Os casos em que a inferência de quantificador é \RightR{\lexists} são ainda
mais fáceis, pois a contração não é necessária e não há condição de variável própria
com que se preocupar. P.ex., suponha que $\RightR{\lexists}$ seja
seguida por $\RightR{\land}$ (desta vez como premissa esquerda). A
sub-!!{proof} relevante é $\pi_1$:
\[
\AxiomC{}
\RightLabel{$\pi_2$}
\Deduce$\Gamma' \fCenter \Delta', \lexists[x][!A(x)], !A(t), !B$
\RightLabel{\RightR{\lexists}}
\UnaryInf$!\Gamma' \fCenter \Delta', \lexists[x][!A(x)], !B$
\AxiomC{}
\RightLabel{$\pi_3$}
\Deduce$\Gamma' \fCenter \Delta', \lexists[x][!A(x)], !C$
\RightLabel{\RightR{\land}}
\BinaryInf$\Gamma' \fCenter \Delta', \lexists[x][!A(x)], !B \land !C$
\DisplayProof
\]
Substitua $\pi_1$ em~$\pi$ por $\pi_1'$:
\[
\AxiomC{}
\RightLabel{$\pi_2$}
\Deduce$\Gamma' \fCenter \Delta', \lexists[x][!A(x)], !A(t), !B$
\AxiomC{}
\RightLabel{$\pi_3'$}
\Deduce$\Gamma' \fCenter \Delta', \lexists[x][!A(x)], !A(t), !C$
\RightLabel{\RightR{\land}}
\BinaryInf$\Gamma' \fCenter \Delta', \lexists[x][!A(x)], !A(t), !B \land !C$
\RightLabel{\RightR{\lexists}}
\UnaryInf$\Gamma' \fCenter \Delta', \lexists[x][!A(x)], !B \land !C$
\DisplayProof
\]
em que $\pi_3'$ é obtida de $\pi_3$ por enfraquecimento com~$!A(t)$ à
direita. Esse enfraquecimento consiste apenas em acrescentar $!A(t)$ ao lado direito
de todo sequente em~$\pi_3$ e, portanto, não altera a ordem de nenhuma
inferência de quantificador. A ordem das inferências de quantificadores é a mesma
em~$\pi'$ que em~$\pi$, exceto pela ordem da inferência \RightR{\lexists}
permutada com~\RightR{\land}, que diminuiu em~$1$. Aplica-se a
hipótese de indução.
\end{proof}

\begin{prob}
Descreva os casos da demonstração de \olref{thm:midsequent} em que
\LeftR{\lexists} termina na premissa de~\RightR{\lif}, e em que
\LeftR{\lforall} termina na premissa esquerda de~\LeftR{\lor}.
\end{prob}

\begin{proof}[Prova do teorema de Herbrand]
Suponha que $\pi$ termine em $\Sequent \lexists[x][!A(x)]$. Por
\olref{thm:midsequent}, podemos transformar $\pi$ em !!a{proof}~$\pi'$
com um sequente médio~$S$, que deve ter a forma \[\Sequent
\lexists[x][!A(x)], !A(t_1), \dots, !A(t_n).\] Como todas as inferências
acima de $S$ são proposicionais, $\lexists[x][!A(x)]$ não pode ser principal
em uma inferência acima de~$S$. Como não é atômica, não pode ser
principal em um axioma. Como $\lexists[x][!A(x)]$ não pode ser uma subfórmula
própria de $!A(t_1)$ (lembre que $!A(x)$ é sem quantificadores), também
não pode estar ativa acima de~$S$. Portanto, apagar $\lexists[x][!A(x)]$ da
sub-!!{proof} que termina em~$S$ produz uma !!{proof} correta
de $\Sequent !A(t_1), \dots, !A(t_n)$, da qual se pode obter !!a{proof} de
$\Sequent !A(t_1) \lor \dots \lor !A(t_n)$ por meio de $n-1$
inferências $\RightR{\lor}$.
\end{proof}

% \begin{prob}
% Find a Herbrand disjunction of $\lexists[x][\lforall[y][(!A(x) \lif
% !A(y))]]$ by finding a cut-free !!{proof} of $\Sequent
% \lexists[x][\lforall[y][(!A(x) \lif !A(y))]]$ and applying the
% transformation in \olref{thm:midsequent} to find a midsequent (if
% necessary).
% \end{prob}

\end{document}
